\scanpage{45}
We shall have in $V$
\[
(4)\qquad\sum_{j\in J}R_j=0
\]
which gives an isomorphism
\[
(5)\qquad V\simeq V_J(\mathbb R)=\operatorname{Ker}(\mathbb R^J\xrightarrow[\text{sum}]{}\mathbb R)
\]
and we have, by the $3^{\text{rd}}$ formula (3)
\[
(6)\qquad\Delta=V^\perp\simeq\mathbb R(\varepsilon)
\]
\margin{``polar axis''}
\margin{orthogonal complement of $V$ in $E$}
\margin{(isomorphism compatible with the quadratic structure)}
thus
\[
(7)\qquad E\simeq V_J(\mathbb R)\oplus\mathbb R(\varepsilon)
\]
\margin{isomorphism of Euclidean spaces}
where in $V_J(\mathbb R)$ we take the unique quadratic form \inserted{$q$} for which the $R_j$ \deleted{are} \inserted{satisfy} \deleted{of square one}
\[
q(R_j)=1
\]
\margin{It is clear, by this characterization of $q$, that the natural action of $G_J$ on $V(J)$ preserves $q$.}
(taking a basis formed of $R_i,R_j$, $R_n$ with coordinates $(-1,-1)$, the quadratic form in question is $x^2-xy+y^2$)

In the Euclidean plane $V$, the $R_i$ \figure{\begin{tikzpicture}[baseline=0pt]\draw[->] (.3,.55)--(0,.35)--(0,0);\end{tikzpicture}} form the vertices of an equilateral triangle inscribed in the unit circle. The action of $G_J$ on $E$ is given by its tautological action \deleted{on} on $V(J)$ and the action by $\operatorname{sg}(g)$ on $\Delta$:
\margin{44}
